\title{Direct Preference Optimization: Your Language Model is Secretly a Reward Model}

\begin{abstract}
Large-scale unsupervised language models (LMs) acquire extensive general knowledge and certain reasoning abilities, but fine-grained modulation of their outputs remains challenging, primarily because their training lacks direct supervision. To address this, prevalent approaches solicit human judgments on the comparative quality of LM outputs and then fine-tune the LMs based on these preference signals—most commonly using reinforcement learning from human feedback (RLHF). However, RLHF involves a multifaceted and frequently unstable process: it initially learns a reward model encapsulating human evaluations, and subsequently applies reinforcement learning to adapt the LM for maximizing this proxy reward, all the while constraining deviations from the original parameters. In this work, we propose an alternative reward model parameterization within RLHF that permits the derivation of the optimal policy analytically in closed form. This reparameterization reframes the RLHF optimization as a standard classification problem, obviating the need for complex reinforcement learning techniques. The approach, which we term \textit{Direct Preference Optimization} (DPO), is marked by its stability, efficiency, and minimal computational demands—removing the necessity for LM sampling during fine-tuning or elaborate hyperparameter adjustment. Experimental evaluations demonstrate that DPO can fine-tune LMs for preference alignment on par with or surpassing conventional methods. In particular, DPO consistently outperforms RLHF with PPO for sentiment control tasks and achieves equivalent or superior performance in summarization and single-turn conversation tasks, all while significantly simplifying the implementation and training pipeline.
\end{abstract}

\section{Introduction}
Large-scale unsupervised language models (LMs) developed on massive corpora demonstrate a range of unexpected competencies~\citep{chowdhery2022palm, brown2020language, touvron2023llama,bubeck2023sparks}. Nevertheless, these models draw their training data from human sources encompassing diverse intentions, priorities, and expertise. Not all of these are appropriate to replicate; for instance, although it is beneficial for an AI code assistant to \textit{recognize} frequent programming errors so it can rectify them, ideally, the model should be skewed towards emulating the (sometimes infrequent) instances of superior coding proficiency that exist within its data. Likewise, it is important that a language model can \textit{identify} widespread misunderstandings, such as a misconception held by half the population, yet it is undesirable for the model to repeat this falsehood as fact half of the time it is asked. In summary, it is vital to distinguish between the model's \emph{preferred outputs and actions} and the broader spectrum of its \textit{acquired information and skills} to construct AI systems that are safe, effective, and amenable to user control~\citep{ouyang2022training}. Although current approaches often rely on reinforcement learning (RL) to align LMs with human preferences, we demonstrate that the RL-inspired learning goal in existing techniques can, in fact, be solved exactly using a straightforward binary cross-entropy loss, thereby streamlining the preference learning workflow.

Broadly speaking, standard practices encode desired conduct in language models through carefully selected collections of human preference data that exemplify behaviors deemed helpful and safe by people. This process of learning preferences is situated after an initial unsupervised pre-training phase on a vast corpus of textual data. The simplest technique for instilling such preferences is through supervised fine-tuning on examples of high-quality human-generated responses. However, the most prevalent and effective strategies employ reinforcement learning from human or AI feedback (RLHF/RLAIF; \citep{christiano2017deep,bai2022constitutional}). In RLHF, a reward model is trained using data of human judgments, which is then used as a signal for RL-based optimization so the language model policy generates outputs that attain high reward while remaining close to its original distribution. Despite the notable capabilities produced by RLHF—especially in conversation and code generation—the associated pipeline is far more intricate compared to supervised methods, involving the training of several LMs, and the necessity to perform policy sampling within the training loop, all of which leads to considerable computational demands.

This work demonstrates a method for optimizing language models to align with human preferences without the need for explicit reward models or reinforcement learning techniques. We introduce \textit{{Direct Preference Optimization} (DPO)}, an approach that implicitly targets the same goal as conventional RLHF algorithms—that is, maximizing a reward subject to a KL-divergence constraint—while being both easy to implement and train. The core idea of the DPO update is to boost the log odds of favored responses relative to less favored ones, incorporating a per-sample importance weighting mechanism. This mechanism mitigates model collapse, a problem observed when employing a straightforward probability ratio objective. DPO, like previous approaches, leverages a theoretical preference model such as the Bradley-Terry model \cite{bradley1952rankanalysis} to assess the alignment between a reward function and actual human preference data. However, unlike standard approaches—which employ the preference model to construct a preference loss used first to fit a reward model and then to train a policy—DPO employs a variable substitution to formulate the preference loss directly as a function of the policy itself. With access to datasets reflecting human choices between model outputs, DPO enables direct policy optimization using a binary cross entropy loss, thus yielding the optimal policy for an implicit reward shaped by human preference.

Our principal contribution is Direct Preference Optimization (DPO): an RL-free, straightforward method for preference-based language model training. Empirical evaluations indicate that DPO performs on par with, or better than, established techniques like PPO-based RLHF across several tasks—including sentiment control, summarization, and conversational agents—on language models containing up to 6B parameters.

\section{Related Work}

As self-supervised language models grow larger, they increasingly demonstrate the ability to perform certain tasks without explicit task-specific training, utilizing zero-shot \citep{radford2019language} or few-shot prompt-based learning \citep{gpt3,megatron,chowdhery2022palm}. Nonetheless, significant gains in downstream task accuracy and adherence to user directives are observed when these models undergo fine-tuning on corpora containing both instructions and human-curated responses \citep{mishra-etal-2022-cross,sanh2022multitask,chung2022scaling,thoppilan2022lamda}. Known as `instruction-tuning', this process not only improves large language models’ capability to handle previously unseen prompts but also broadens their practical utility \citep{chung2022scaling}. While instruction-tuning has proven effective, it is often more feasible to gather human preference data in the form of comparative judgments, as opposed to accumulating high-quality expert-written samples. As a result, later research efforts have adapted LLMs using collections of pairwise human preferences, leading to enhancements in machine translation \citep{kreutzer-etal-2018-reliability}, summarization \citep{stiennon2022learning,ziegler2020finetuning}, narrative generation \citep{ziegler2020finetuning}, and response to instructions \citep{ouyang2022training,ramamurthy2023is}. Such approaches typically begin by learning a neural reward model aligned to these preferences—using frameworks like the Bradley-Terry model \citep{bradley1952rankanalysis}—followed by language model fine-tuning via reinforcement learning algorithms, such as REINFORCE \citep{williams1992reinforce}, proximal policy optimization (PPO; \cite{schulman2017proximal}), or related strategies \citep{ramamurthy2023is}. Related research lines employ LLMs, themselves tuned with human feedback, to generate additional synthetic preference annotations for specific attributes, including safety and non-harmfulness \citep{bai2022constitutional}; here, limited human involvement is required, often restricted to providing textual guidelines for model annotation. Collectively, these advancements bridge earlier efforts: training language models through reinforcement learning toward diverse goals~\citep{Ranzato2015SequenceLT,paulus2018a,wu2018learning} and learning directly from human preferences \citep{christiano2017deep,kupcsik2018learning}. Although leveraging relative human judgments presents an attractive methodology, scaling reinforcement learning-based fine-tuning for large language models is fraught with practical difficulties; the present work introduces a theoretically-motivated method for optimizing with relative preferences that circumvents the need for reinforcement learning.

Beyond linguistic applications, policy learning from preferences has been explored in both the contextual bandit and reinforcement learning frameworks, resulting in a variety of proposed methodologies. In the contextual bandit literature, learning from preferences or action rankings rather than from direct reward feedback leads to what is called the contextual dueling bandit (CDB; \cite{yue2012karmed,dudik2015contextual}). In scenarios where explicit rewards are unavailable, theoretical work on CDBs redefines the goal of optimal policy in terms of identifying a \textit{von Neumann winner}: a policy guaranteed to outperform \textit{any} competing policy at least half the time in expectation \citep{dudik2015contextual}. Notably, the CDB setting relies on online provision of preference signals, in contrast to typical preference learning from humans, where a static, finite set of preference-annotated action pairs is available for offline use \citep{yan2022human}. Similarly, \textit{preference-based RL} (PbRL) utilizes binary preference comparisons produced by an underlying, hidden `scoring' function, as opposed to numeric rewards \citep{BusaFekete2014,ruiz2023dueling}. Although a variety of algorithms exist for PbRL—including approaches that support leveraging previously collected off-policy preference data—these methods generally involve a two-step process: first estimating the latent scoring (reward) function, and then optimizing policy with respect to it \citep{jain2013learning,BusaFekete2014,christiano2017deep,sadigh2017active,kupcsik2018learning}. In contrast, we propose a unified, direct policy optimization framework that trains a policy to adhere to given preferences in a single stage.

\section{Preliminaries}\label{section:prelims}

We summarize the RLHF pipeline as outlined by \citeauthor{ziegler2020finetuning} and expanded upon in later works (\citep{stiennon2022learning, bai2022training, ouyang2022training}). This approach commonly involves three key steps: (1) supervised fine-tuning (SFT); (2) acquisition of preference data and construction of a reward model; and (3) reinforcement learning optimization.

\textbf{SFT}: RLHF procedures begin with supervised fine-tuning of a pre-trained language model on carefully curated, high-quality datasets relevant to the downstream tasks (such as dialogue systems, summarization, etc.), resulting in a model $\pi^\text{SFT}$.

\textbf{Reward Modelling Phase}: In the next stage, prompts $x$ are input to the SFT model to generate answer pairs $(y_1, y_2)\sim \pi^\text{SFT}(y \mid x)$. These answer pairs are then assessed by human annotators, who indicate a preference (notated as $y_w\succ y_l \mid x$, where $y_w$ and $y_l$ are the more and less favored responses, respectively). These preferences are modeled as being produced according to an unknown latent reward function $r^*(y, x)$, which is not directly observable. To represent such preferences, various probabilistic models are available; among them, the Bradley-Terry (BT) \cite{bradley1952rankanalysis} model is frequently used, although the more general Plackett-Luce ranking models \citep{plackett1975analysis, luce2012individual} can also be applied when multiple ranked answers are collected. Under the BT model, the probability distribution over human preferences, denoted as $p^*$, is given by:

\begin{equation}\label{eq:bradley-terry}
    p^*(y_1\succ y_2 \mid x)=\frac{\exp\left(r^*(x, y_1)\right)}{\exp\left(r^*(x, y_1)\right) + \exp\left(r^*(x, y_2)\right)}.
\end{equation}

Given a fixed dataset of preference comparisons, $\mathcal{D}=\bigl\{x^{(i)}, y_w^{(i)}, y_l^{(i)}\bigr\}_{i=1}^N$, which are drawn according to $p^*$, one can represent the reward with a parameterized model $r_{\phi}(x, y)$ and optimize its parameters using maximum likelihood estimation. When interpreting the comparison task as a binary classification problem, the corresponding negative log-likelihood loss takes the form:

\begin{equation}\label{eq:reward_model}
    \mathcal{L}_R(r_{\phi}, \mathcal{D}) = -\mathbb{E}_{(x, y_w, y_l)\sim \mathcal{D}}\bigl[\log \sigma(r_{\phi}(x, y_w)- r_{\phi}(x, y_l))\bigr]
\end{equation}

where $\sigma$ denotes the logistic sigmoid function. For large language models, the reward model $r_{\phi}(x, y)$ is often seeded with parameters from the SFT model $\pi^\text{SFT}(y \mid x)$, to which a linear layer is appended atop the final transformer output to yield a single scalar reward prediction \cite{ziegler2020finetuning}. To reduce the variability in the reward function, previous works standardize the predicted reward so that $\mathbb{E}_{x,y\sim \mathcal{D}}\left[r_\phi(x, y)\right] = 0$ for every $x$.

\textbf{RL Fine-Tuning Phase}: During the reinforcement learning stage, this trained reward model serves as the feedback mechanism to update the language model. Following earlier studies~\citep{jaques2017sequence, jaques2020human}, the objective is framed as

\begin{equation}\label{eq:RL}
\max_{\pi_{\theta}}  \mathbb{E}_{x\sim \mathcal{D}, y\sim \pi_{\theta}(y \mid x)}\bigl[r_{\phi}(x, y)\bigr] - \beta\mathbb{D}_{\textrm{KL}}\bigl[\pi_{\theta}(y\mid x)\mid \mid \pi_\text{ref}(y\mid x)\bigr],
\end{equation}

where $\beta$ adjusts the allowable divergence from a reference policy $\pi_\text{ref}$, typically taken as the SFT baseline $\pi^\text{SFT}$. 
Commonly, $\pi_\theta$ is also initialized using $\pi^\text{SFT}$. This regularization is crucial, as it ensures that the model remains close to the regime where the reward estimator is reliable, while also preserving diversity of responses and avoiding concentration on a single output mode. Because text generation is inherently discrete and the objective is not differentiable, this problem is usually addressed via reinforcement learning. In particular, prior works~\citep{ziegler2020finetuning, stiennon2022learning, bai2022training, ouyang2022training} recast the reward as ${r(x, y) = r_{\phi}(x, y) -\beta (\log \pi_{\theta}(y\mid x) - \log \pi_\text{ref}(y\mid x))}$, and apply the PPO algorithm for policy optimization \cite{schulman2017proximal}. 

\section{Direct Preference Optimization}\label{sec:DPO}

To overcome the scalability bottlenecks associated with RL-based fine-tuning for language models, we propose a streamlined method that directly optimizes policies using preference data. Departing from conventional RLHF pipelines—which separately learn a reward function before optimizing it via reinforcement learning—our framework exploits a particular parameterization of the reward such that the resulting optimal policy can be determined analytically, thereby eliminating the need for iterative RL updates. As detailed below, the core idea is to use a closed-form mapping from the space of rewards to the space of optimal policies, enabling reformulation of reward-based losses into policy-based ones. Through this change-of-variables mechanism, we bypass the need to explicitly fit a separate reward network, yet still optimize using established human preference models like Bradley-Terry. In this construction, the policy model encapsulates both language generation and

\textbf{Deriving the DPO objective.} We begin with the reinforcement learning (RL) objective from previous literature, specifically Eq.~\ref{eq:RL}, under an arbitrary reward function $r$. Consistent with prior studies~\citep{peters2007reinforcement, peng2019advantage, korbak2022reinforcement, go2023aligning}, it can be readily established that the solution optimizing the KL-regularized reward objective in Eq.~\ref{eq:RL} is given by:

\begin{equation}\label{eq:op_policy}
    \pi_r(y\mid x) = \frac{1}{Z(x)}\pi_\text{ref}(y\mid x)\exp\left(\frac{1}{\beta}r(x, y)\right),
\end{equation}

with the partition function $Z(x) =\sum_{y}\pi_\text{ref}(y\mid x)\exp\left(\frac{1}{\beta}r(x, y)\right)$. Full details are provided in Appendix~\ref{app:derivation1}. However, even when substituting an MLE-based estimate $r_{\phi}$ for the true reward $r^*$, evaluating $Z(x)$ remains computationally intensive~\citep{korbak2022reinforcement, go2023aligning}, limiting practical applicability. To circumvent this issue, Eq.~\ref{eq:op_policy} can be reformulated to express the reward function using the optimal policy $\pi_r$, the reference policy $\pi_\text{ref}$, and the (unknown) partition function $Z(\cdot)$. Specifically, by applying a logarithm to both sides of Eq.~\ref{eq:op_policy} and simplifying, we get:

\begin{equation}\label{eq:main_eq}
    r(x,y) =\beta \log \frac{\pi_r(y\mid x)}{\pi_\text{ref}(y\mid x)} + \beta \log Z(x).
\end{equation}

This parameterization can likewise be applied to the actual reward $r^*$ and the associated optimal policy $\pi^*$. The Bradley-Terry model, crucially, only involves the reward differences between output candidates, specifically, ${p^*(y_1 \succ y_2 \mid x) = \sigma(r^*(x, y_1) - r^*(x, y_2))}$. When substituting the form for $r^*(x,y)$ from Eq.~\ref{eq:main_eq} into Eq.~\ref{eq:bradley-terry}, the partition function terms negate each other, enabling us to represent the human preference probability exclusively in terms of $\pi^*$ and $\pi_\text{ref}$. Thus, the optimal RLHF policy $\pi^*$ in the Bradley-Terry framework satisfies the following condition:

\begin{equation}\label{eq:objective}
    p^*(y_1\succ y_2 \mid x)=\frac{1}{1 + \exp\left(\beta \log \frac{\pi^*(y_2\mid x)}{\pi_\text{ref}(y_2\mid x)} - \beta \log \frac{\pi^*(y_1\mid x)}{\pi_\text{ref}(y_1\mid x)}\right)}
\end{equation}

See Appendix~\ref{app:derivation2} for a full derivation. While Eq.~\ref{eq:objective} is expressed for the Bradley-Terry model, an analogous derivation can be conducted for the broader Plackett-Luce family~\citep{plackett1975analysis, luce2012individual}, as provided in Appendix~\ref{app:plackett_luce_models}.

Given this expression for the likelihood of human preference data in terms of the optimal policy, it becomes possible to frame a maximum likelihood learning objective for a parameterized policy $\pi_\theta$. Similar to the procedure in reward modeling (cf. Eq.~\ref{eq:reward_model}), this leads us to define the policy objective as:

\begin{equation}\label{eq:optimum_model}
    \mathcal{L}_\text{DPO}(\pi_{\theta}; \pi_\text{ref}) = -\mathbb{E}_{(x, y_w, y_l)\sim \mathcal{D}}\left[\log \sigma \left(\beta \log \frac{\pi_{\theta}(y_w\mid x)}{\pi_\text{ref}(y_w\mid x)} - \beta \log \frac{\pi_{\theta}(y_l\mid x)}{\pi_\text{ref}(y_l\mid x)}\right)\right].
\end{equation}

In this approach, we model the implicit reward through an alternative parameterization such that its optimal policy directly corresponds to $\pi_\theta$. Additionally, our methodology aligns with fitting a reparameterized Bradley-Terry model, thereby inheriting desirable theoretical guarantees, including consistency under appropriate distributional assumptions on the preference data \cite{bong2022generalized}. Further discussion of the theoretical implications of DPO, especially in comparison to related frameworks, is provided in Section~\ref{sec:theory}.

\textbf{Mechanistic insights into the DPO update.} To elucidate the operational mechanics of DPO, examining the gradient of its loss function, $\mathcal{L}_\text{DPO}$, proves instructive. The gradient with respect to $\theta$ is given by:

\begin{multline*}\label{eq:gradient}
    \nabla_\theta \mathcal{L}_\text{DPO}(\pi_\theta;\pi_\text{ref}) = \\ -\beta\mathbb{E}_{(x, y_w, y_l) \sim \mathcal{D}} \bigg[\underbrace{\sigma(\hat{r}_\theta(x, y_l) - \hat{r}_\theta (x, y_w))}_\text{higher weight when reward estimate is wrong}\bigg[\underbrace{\nabla_\theta\log \pi(y_w \mid x)}_\text{increase likelihood of $y_w$} - \underbrace{\nabla_\theta\log\pi(y_l \mid x)}_\text{decrease likelihood of $y_l$}\bigg]\bigg],
\end{multline*}

where the implicit reward associated with each completion, $\hat{r}_\theta(x, y)$, is computed as $\beta \log \frac{\pi_\theta(y \mid x)}{\pi_\text{ref}(y \mid x)}$, leveraging both the current model $\pi_\theta$ and a reference model $\pi_\text{ref}$ (see Section~\ref{sec:theory} for additional context). The intuition is that the loss gradient updates the model to elevate the probability assigned to preferred completions $y_w$ and suppress that of less-preferred completions $y_l$. Critically, the update places greater emphasis on instances where the implicit reward $\hat{r}_\theta$ overvalues the dispreferred response, scaled by $\beta$. Thus, the update magnitude reflects the extent of misordering by the reward function, incorporating the impact of the KL regularization. Empirical evidence from our experiments highlights the necessity of this weighting mechanism, as omitting the weighting factor may lead to pathological behavior in the model (see Appendix Table~\ref{tab:unlikelihood_generations}).

\textbf{DPO overview.} The standard DPO workflow involves the following steps: 1) For each prompt $x$, draw completion samples $y_1, y_2 \sim \pi_\text{ref}(\cdot \mid x)$, obtain human preference annotations to assemble an offline preferences dataset $\mathcal{D} = \{x^{(i)}, y_w^{(i)}, y_l^{(i)}\}_{i=1}^N$; 2) update the language model $\pi_\theta$ by minimizing $\mathcal{L}_\text{DPO}$, conditioning on the chosen $\pi_\text{ref}$, dataset $\mathcal{D}$, and a specified $\beta$. 

In real-world applications, it is often preferable to make use of publicly available preference datasets rather than producing new samples and collecting fresh human annotations. Typically, since these datasets originate from $\pi^\text{SFT}$, we use $\pi_\text{ref} = \pi^\text{SFT}$ when possible. In cases where $\pi^\text{SFT}$ is not accessible, we set $\pi_\text{ref}$ by maximizing the likelihood of the preferred completions ${(x, y_w)}$, i.e., we let ${\pi_\text{ref} = \argmax_{\pi}\mathbb{E}_{x, y_w \sim \mathcal{D}}\left[\log \pi(y_w \mid x)\right]}$. This initialization is designed to reduce the mismatch between the unknown true reference distribution and the one employed by DPO, $\pi_\text{ref}$. Additional implementation details and selected hyperparameters are presented in Appendix~\ref{app:implementation}.

\section{Theoretical Analysis of DPO} In this part, we further interpret the workings of DPO, present theoretical foundations, and discuss how DPO resolves challenges found in actor-critic methods for RLHF—such as those encountered with PPO~\cite{schulman2017proximal}.

\label{sec:theory}

\subsection{Your Language Model Is Secretly a Reward Model} DPO uniquely sidesteps the need for explicit reward modeling and reinforcement learning in policy training by employing a single maximum likelihood objective. Specifically, the objective in Eq. \ref{eq:main_eq} is equivalent to using the Bradley-Terry model, where the reward is parameterized as $r^*(x, y) = \beta \log\frac{\pi^*_\theta(y \mid x)}{\pi_\text{ref}(y \mid x)}$; here, the parametric model $\pi_{\theta}$ is trained analogously to optimizing a reward model, as in Eq. \ref{eq:reward_model}, following a change of variables. The subsequent discussion develops the underlying theory supporting this reparameterization, demonstrates that it imposes no restrictions on the set of possible learned reward models, and proves that it enables precise recovery of the optimal policy. Our starting point is the establishment of an equivalence relation between reward functions.

\begin{definition}
Two reward functions $r(x, y)$ and $r'(x, y)$ are said to be equivalent if and only if ${r(x, y)-r'(x, y) = f(x)}$ for some function $f$.
\end{definition}

One can readily verify that this definition establishes an equivalence relation, thereby partitioning the collection of reward functions into equivalence classes. This leads to the following two results:

\begin{lemma}\label{lemma:same_prefrence}
Within the Plackett-Luce, and specifically the Bradley-Terry, framework for preferences, any two reward functions belonging to the same equivalence class yield identical distributions over preferences.
\end{lemma}

\begin{lemma}\label{lemma:same_policy}
    For the constrained RL setting, all reward functions that are members of the same equivalence class correspond to the same optimal policy.
\end{lemma}

The arguments supporting these lemmas are elementary and can be found in Appendix \ref{app:lemma1}. The first lemma highlights the well-documented issue of under-specification associated with the Plackett-Luce family of models \cite{plackett1975analysis}. To address this ambiguity, additional identifiability conditions are usually required to obtain meaningful MLE estimates for Eq. \ref{eq:reward_model} \cite{bong2022generalized}. The second lemma indicates that all reward functions within an equivalence class lead to an identical optimal policy; therefore, our main objective is to recover any representative from the correct class of reward functions. The following theorem, which is proven in Appendix~\ref{app:thm1}, formalizes this:

\begin{theorem}\label{thm:main}
    Given mild conditions, every reward class that is compatible with the Plackett-Luce (and, in particular, the Bradley-Terry) models admits a representation of the form ${r(x, y) = \beta \log \frac{\pi(y\mid x)}{\pi_\text{ref}(y\mid x)}}$, for some model $\pi(y\mid x)$ and a specified reference model $\pi_\text{ref}(y \mid x)$.
\end{theorem}

\begin{sproof}
    Take any reward function $r(x, y)$, and let $\pi_r(y \mid x)$ denote its associated optimal model as given by Eq. \ref{eq:op_policy}. We aim to show that an element of the equivalence class containing $r$ can be expressed using the reparameterization above. Define the transformation $f$ as
\begin{equation}
    f(r; \pi_\text{ref}, \beta)(x, y) = r(x, y) - \beta\log\sum_{y}\pi_\text{ref}(y\mid x)\exp\left(\frac{1}{\beta}r(x, y)\right)
\end{equation}
Here, $f$ adjusts the original reward function by subtracting the logarithm of the corresponding partition function, scaled by $\beta$. Because this normalization term is dependent solely on $x$, $f(r; \pi_\text{ref}, \beta)(x, y)$ remains in the same equivalence class as $r(x, y)$. Substituting $r$ by the right side of Eq.~\ref{eq:main_eq}, which is applicable to any reward function, results in $f(r; \pi_\text{ref}, \beta)(x, y) = \beta \log \frac{\pi_r(y\mid x)}{\pi_\text{ref}(y\mid x)}$. Thus, the projection $f$ yields a member of the equivalence class with the required structure, ensuring the proposed reparameterization is sufficiently general for our reward modeling.
\end{sproof}

An alternative perspective on Theorem~\ref{thm:main} is that it delineates the specific reward function, among all those in the same equivalence class, selected by the DPO reparameterization—namely, the reward function fulfilling

\begin{equation}\label{eq:lag_p}
     \sum_{y}\underbrace{\pi_\text{ref}(y\mid x)\exp\left(\frac{1}{\beta}r(x, y)\right)}_{=\pi(y\mid x)\text{, by Thm.~\ref{thm:main} reparam.}} = 1,
\end{equation}

In other words, $\pi(y\mid x)$ constitutes a valid probability distribution, as its entries are nonnegative and sum to one. Additionally, by considering Eq.~\ref{eq:op_policy}, Eq.~\ref{eq:lag_p} represents the partition function associated with the optimal policy derived from the reward function $r(x, y)$. The principal idea underlying the DPO algorithm is to introduce specific constraints into the otherwise underdetermined Plackett-Luce (and in particular, Bradley-Terry) class of preference models. By doing so, one retains the expressive capacity of the representable reward model class, while ensuring that the optimal policy in Eq.~\ref{eq:op_policy} is analytically tractable for arbitrary prompts $x$.

\subsection{Instability of Actor-Critic Algorithms}
Our framework also facilitates the analysis of instability phenomena characteristic of standard actor-critic approaches in RLHF, such as PPO. Focusing on the RL fine-tuning phase discussed in Section \ref{section:prelims}, we can relate this analysis to the control as inference perspective \cite{levine2018reinforcement} as applied to the constrained RL formulation in \ref{eq:RL}. Suppose we employ a parameterized policy $\pi_{\theta}(y\mid x)$, and seek to minimize $\mathbb{D}_{\text{KL}}[\pi_{\theta}(y|x) \mid \mid \pi^*(y\mid x)]$, with $\pi^*$ denoting the optimal policy according to Eq.~\ref{eq:optimum_model} and induced by reward function $r_{\phi}(y, x)$. Some algebraic manipulation leads us to the following optimization problem:

\begin{equation}\label{eq:AC}
    \max_{\pi_{\theta}}\mathbb{E}_{\pi_{\theta}(y\mid x)}\bigg[\underbrace{r_{\phi}(x, y) -\beta\log\sum_{y}\pi_\text{ref}(y\mid x)\exp\left(\frac{1}{\beta}r_{\phi}(x, y)\right)}_{f(r_{\phi}, \pi_\text{ref}, \beta)} - \underbrace{\beta\log\frac{\pi_{\theta}(y\mid x)}{\pi_\text{ref}(y\mid x)}}_{\text{KL}}\bigg]
\end{equation}

The objective considered here aligns with those optimized in earlier studies \citep{ziegler2020finetuning, stiennon2022learning, bai2022training, ouyang2022training}, utilizing the DPO-equivalent reward formulation for the reward class $r_{\phi}$. Within this framework, the normalization component in $f(r_{\phi}, \pi_\text{ref}, \beta)$ can be viewed as the soft value function associated with the reference policy $\pi_\text{ref}$. Although this normalization factor does not influence the solution at optimality, omitting it can result in high-variance policy gradients, potentially causing instability during learning. One way to address this is by introducing a learned value function to account for normalization; however, optimizing such a value function is often challenging. Prior literature instead normalizes rewards by referencing a human completion baseline, which essentially provides a single-sample Monte-Carlo estimate for the normalization constant. By contrast, the DPO reparameterization offers a reward formulation that does not necessitate any baseline.

\section{Experiments}
This section empirically assesses DPO’s capability to train policies directly from preference feedback. We begin with a controlled text generation scenario to investigate how effectively DPO manages the tradeoff between reward maximization and minimization of KL-divergence from the reference policy, relative to standard preference learning methods like PPO. Subsequently, we test DPO on larger-scale models and more challenging RLHF tasks, including dialogue and summarization. Our results indicate that, with minimal hyperparameter tuning, DPO generally matches or surpasses established baselines such as RLHF with PPO, as well as the strategy of selecting the highest-scoring trajectory from $N$ samples using a learned reward function. We detail our experimental methodology prior to presenting results, with further specifics available in Appendix~\ref{app:exp_details}.

\textbf{Tasks.} Our study investigates three distinct open-domain text generation tasks. In each scenario, the methods under consideration learn a policy from a dataset of preferences denoted as $\mathcal{D}=\bigl\{x^{(i)}, y_w^{(i)}, y_l^{(i)}\bigr\}_{i=1}^N$. For \textbf{controlled sentiment generation}, the input $x$ is a partial movie review selected from the IMDb dataset \cite{maas-EtAl:2011:ACL-HLT2011}, and the objective for the policy is to complete the review with text $y$ expressing positive sentiment. To facilitate controlled comparison, we synthetically create preference pairs between generated outputs by utilizing a sentiment classifier, ensuring that $p(\text{positive}\mid x,y_w)>p(\text{positive}\mid x,y_l)$. The SFT model is developed by fine-tuning GPT-2-large to convergence using review data from the IMDb training partition (see further explanation in App~\ref{app:sentiment_details}). For the \textbf{summarization} task, $x$ consists of Reddit forum posts, where the policy aims to generate a summary $y$ that encapsulates the main points. In line with earlier research, our experiments utilize the Reddit TL;DR summarization corpus \citep{volske-etal-2017-tl}, augmented with human preference data from \citeauthor{stiennon2022learning}. Here, the SFT model is adapted using human-written summaries of forum posts,\footnote{\url{https://huggingface.co/CarperAI/openai_summarize_tldr_sft}} and RLHF is implemented with the TRLX framework \citep{leandro_von_werra_2023_7790115}. The preference data used comes from human annotations of outputs from a distinct, though similarly fine-tuned, SFT system described by \citeauthor{stiennon2022learning}. Lastly, the \textbf{single-turn dialogue} setting takes $x$ to be a user question—which can span diverse topics such as science or personal advice—with the policy's role being to deliver an informative and appropriate response $y$. For this, we employ the Anthropic Helpful and Harmless dialogue dataset \citep{bai2022training}, featuring 170,000 dialogues between human users and an AI assistant. Each entry concludes with two candidate responses generated by a (large but unidentified) language model, accompanied by a human-annotated preference for one of the replies. In this scenario, there is no pre-existing SFT model; instead, we construct one by fine-tuning a publicly available language model on the subset of responses identified as preferred by annotators.

\textbf{Evaluation.} Our experimental assessment employs two distinct strategies. To investigate how effectively each algorithm addresses the constrained reward maximization problem, we consider the controlled sentiment generation scenario, where performance is measured by plotting the set of achievable trade-offs between reward and KL-divergence relative to the reference policy; here, this is feasible because the true reward function (given by a sentiment classifier) is accessible. In contrast, under more realistic conditions where the reward function is not explicitly available, we rely on a \textit{win rate} metric relative to a baseline policy, enlisting GPT-4 to approximate human judgment in evaluating summary quality and helpfulness of dialogue responses. For the summarization task, the baseline comprises the test set’s reference summaries; for dialogue, the baseline corresponds to the preferred reply from the test split. While literature has reported that language models can often serve as superior automated evaluators compared to standard metrics \citep{Chen2023ExploringTU}, we supplement this with a human evaluation in Sec.~\ref{sec:human-judgments} to substantiate our use of GPT-4 for evaluation. Our results show a strong alignment between GPT-4 and human raters, with the level of agreement between humans and GPT-4 meeting or exceeding typical inter-annotator agreement among humans.

\textbf{Methods.} Alongside DPO, we include several other established strategies for training language models in line with human preferences. For summarization, we evaluate zero-shot prompting using \textbf{GPT-J} \citep{gpt-j}, while the dialogue task employs 2-shot prompting with \textbf{Pythia-2.8B} \citep{biderman2023pythia}. Additional baselines involve the \textbf{SFT} model and \textbf{Preferred-FT}, the latter obtained by supervised fine-tuning on the preferred outputs $y_w$, using either the SFT model for sentiment and summarization, or a generic language model for single-turn dialogue. We also incorporate \textbf{Unlikelihood}~\citep{welleck2019neural}, a pseudo-supervised strategy that enhances the probability of $y_w$ while actively decreasing that of $y_l$, modulated by an optional scaling factor $\alpha\in[0,1]$ for the unlikelihood component. Additionally, our analysis covers \textbf{PPO} \citep{schulman2017proximal}, which employs a reward function inferred from preference annotations, as well as \textbf{PPO-GT}—an oracle variant utilizing access to the ground-truth reward function, which is restricted to the controlled sentiment setting. For sentiment experiments, two PPO-GT implementations are compared: an out-of-the-box solution \cite{leandro_von_werra_2023_7790115} and an enhanced version featuring reward normalization and hyperparameter refinements (these same improvements are applied when evaluating standard PPO with learned rewards). We also explore the \textbf{Best of $N$} approach, which entails generating $N$ outputs from the SFT model (or from Preferred-FT for dialogue) and returning the response with the highest score per the learned reward model. While this technique can yield strong performance by isolating the reward model quality from PPO training, its reliance on sampling $N$ completions per test query renders it computationally expensive for all but very small values of $N$.

\subsection{How well can DPO optimize the RLHF objective?}

The reward maximization objective used in standard RLHF methods involves a KL-divergence constraint, effectively trading off between maximizing reward and keeping the policy close to a designated reference policy. Thus, any comparison between methods must jointly consider the magnitude of reward and the corresponding KL-divergence; gaining marginally more reward at the cost of significantly increased KL may not be advantageous. Figure~\ref{fig:frontier-tldr-main} presents the reward-KL tradeoff curves for several algorithms within the sentiment task. For each algorithm, we perform several training sessions, varying a key hyperparameter related to policy conservativeness (target KL $\in\{3,6,9,12\}$ for PPO, $\beta \in \{0.05,0.1,1,5\}$, $\alpha\in\{0.05,0.1,0.5,1\}$ for unlikelihood, and different random seeds for preferred-FT), leading to a total of 22 experimental runs. After every 100 training iterations, until convergence, we evaluate each trained policy on a fixed set of test prompts, measuring the mean reward under the true reward function and also the mean sequence-level KL\footnote{Specifically, the aggregate KL-divergence over all time steps.} relative to the reference policy, denoted $\text{KL}\left(\pi\mid \mid \pi_\text{ref}\right)$. Our results show that DPO produces a significantly more favorable reward-KL frontier than competing methods, attaining the highest rewards for comparably low KL values. This finding is especially striking for two main reasons. First, despite DPO and PPO both directly optimizing the same underlying objective, DPO is much more effective in practice, with its reward-KL curve strictly outperforming that of PPO. Second, DPO yields a superior tradeoff even in comparison to PPO equipped with access to the true reward values (PPO-GT).

\subsection{Can DPO scale to real preference datasets?}
\label{sec:dpo-real-datasets}
We proceed by assessing the fine-tuning capabilities of DPO in the contexts of summarization and single-turn dialogue tasks. In summarization, automated metrics such as ROUGE are often not well aligned with human judgments~\citep{stiennon2022learning}, and existing research has shown that applying PPO-based fine-tuning on human preference data can yield superior summaries. Our evaluation consists of generating model completions from the TL;DR summarization dataset's test split and calculating the mean win rate compared to reference completions. For all approaches, outputs are sampled across a temperature range from 0.0 to 1.0, and their win rates are depicted in Figure~\ref{fig:frontier-tldr-main} (right). The models DPO, PPO, and Preferred-FT are all fine-tuned beginning with the same GPT-J SFT checkpoint\footnote{\url{https://huggingface.co/CarperAI/openai_summarize_tldr_sft}}. The results indicate that DPO attains a win rate near 61\% at a temperature of 0.0, surpassing PPO's performance, which reaches about 57\% at its optimal sampling temperature. Additionally, DPO secures a higher maximum win rate than the best-of-$N$ baseline. It is important to highlight that we did not thoroughly optimize DPO's $\beta$ hyperparameter, suggesting the reported results could be a conservative estimate of its true effectiveness. DPO further demonstrates greater stability across temperature variations compared to PPO, whose performance can decline toward that of the baseline GPT-J model when sampled at higher temperatures. Preferred-FT, in contrast, does not exhibit significant improvements over the SFT baseline. We also conduct direct human comparison studies between DPO and PPO in Section~\ref{sec:human-judgments}, where samples generated by DPO at a temperature of 0.25 were favored 58\% of the time over PPO outputs sampled at 0.

For single-turn dialogue, we assess various approaches on the subset of the test split from the Anthropic HH dataset \citep{bai2022training}, focusing specifically on instances involving a single exchange between human and assistant. In our GPT-4-based assessments, preferred test completions serve as ground truth to calculate the win rate of each method. Since a standard SFT model is not established for this task, we initialize from the pre-trained Pythia-2.8B, apply Preferred-FT on the selected completions to create a reference model whose outputs align with the completion distribution, and subsequently train using DPO. For comparison, we include the strongest outcome among 128 Preferred-FT completions—determined after observing that increasing $N$ in the Best of $N$ baseline yields diminishing returns at $N=128$ for this setting (see Appendix Figure~\ref{fig:best-of-n})—as well as a 2-shot prompted baseline with Pythia-2.8B. Our experiments indicate that DPO achieves comparable or superior results relative to alternatives when using the optimal sampling temperature for each. Additionally, we analyze the performance of an RLHF agent trained via PPO on the Anthropic HH dataset \footnote{\url{https://huggingface.co/reciprocate/ppo_hh_pythia-6B}} as released by a reputable source \footnote{\url{https://github.com/CarperAI/trlx/tree/main/examples/hh}}, but are unable to identify any prompting or temperature configuration that surpasses the baseline Pythia-2.8B model’s performance. Drawing on insights from TL;DR and the alignment in reward function optimization, we treat Best of 128 as a rough surrogate for PPO performance. In summary, DPO stands out as the only method that both enhances over preferred completions within the Anthropic HH dataset and matches or exceeds the more resource-intensive Best of 128 approach, all while maintaining computational efficiency. Moreover, as depicted in Figure~\ref{fig:dialogue-main}, DPO attains its optimal results with relatively rapid convergence.

\subsection{Generalization to a new input distribution}

To further assess how PPO and DPO cope with distributional changes, we analyze the policies learned from our Reddit TL;DR summarization experiment when deployed on an out-of-distribution setting: the test split of the CNN/DailyMail dataset \citep{nallapati-etal-2016-abstractive}, which consists of news articles. For this evaluation, we utilize the optimal sampling temperatures identified on TL;DR (0 and 0.25). The findings, shown in Table~\ref{tab:ood}, are based on win rates against the reference summaries as assessed by GPT-4, employing the same prompt as for TL;DR, except for substituting “forum post” with “news article.” On this novel input distribution, DPO consistently surpasses PPO by a notable margin. These results serve as preliminary support that DPO can match or exceed the generalization performance of PPO policies, despite DPO not utilizing the supplementary unlabeled Reddit TL;DR prompts incorporated by PPO.

\subsection{Validating GPT-4 judgments with human judgments}
\label{sec:human-judgments}
To examine the validity of GPT-4’s evaluations, we conduct a human rating study drawing from the TL;DR summarization experiment and two variants of the GPT-4 prompt. The \textbf{GPT-4 (S)} (simple) prompt requests a decision as to which summary more effectively captures the important content of the post, whereas the \textbf{GPT-4 (C)} (concise) prompt asks the same but also considers conciseness—an adjustment made since GPT-4, under the simple prompt, tends to prefer summaries that are longer and more repetitive than those favored by human annotators. Full prompt texts are included in Appendix~\ref{app:prompts}. Our comparisons involve three systems representing a range of output quality: the top-performing (DPO, temp. 0.25), bottom-performing (PPO, temp. 1.0), and an intermediate (SFT, temp. 0.25) method, each paired against PPO using its best-performing temperature (greedy sampling). In both prompting conditions, we observe that GPT-4’s preferences align with human preferences at a rate comparable to inter-human agreement (with multiple human raters available only for the DPO and PPO-1 matchups due to rater limitations). Notably, the \textbf{GPT-4 (C)} prompt produces win rates more in line with human ratings, which motivates its use as our standard for the principal results in Section~\ref{sec:dpo-real-datasets}. Further methodological details—including the rating interface and the list of participants—can be found in Appendix~\ref{app:human-study}.

\section{Discussion}
Preference-based learning stands out as a robust and adaptable approach for developing proficient and aligned language models. In this work, we presented DPO, a straightforward methodology that allows training language models based on preferences, entirely bypassing reinforcement learning. Instead of reconfiguring the preference learning framework to align with conventional RL formulations for the sake of utilizing standard RL algorithms, DPO establishes a correspondence between language model policies and reward functions, facilitating the direct optimization of language models to reflect human preferences via a basic cross-entropy loss. This approach neither requires reinforcement learning nor sacrifices generality. Even with minimal hyperparameter adjustment, DPO achieves results that are on par with, or superior to, prominent RLHF methods, such as those utilizing PPO, thereby significantly lowering the entry barrier for preference-based language model training.

\textbf{Limitations \& Future Work.} Several critical questions arise from our findings that warrant future investigation. One is the extent to which DPO-trained policies maintain their performance on out-of-distribution tasks compared to those trained with explicit reward models. Preliminary evidence indicates DPO's generalization capabilities are comparable to PPO-based methods, though more thorough evaluation is necessary. An additional direction includes assessing the effectiveness of leveraging unlabeled data via self-labeling from the DPO policy. It remains an open question how reward over-optimization appears within the direct preference optimization framework, and whether the modest performance decline observed in Figure~\ref{fig:dialogue-main}-right is symptomatic of this phenomenon. Our experiments were limited to models containing up to 6B parameters; thus, investigating the scalability of DPO to much larger, state-of-the-art models represents a promising research avenue. In terms of evaluation, we observed that win rates determined by GPT-4 can depend on the prompt formulation; refining automated evaluation methods to elicit reliable judgments is another pertinent subject for further work. Lastly, DPO's potential applications are broad and extend beyond language modeling based on human preferences, offering opportunities for advancing preference-based training across various generative modeling domains.